%%
%% SIT330-770 Distinction Task 2.2D — Introduction (skeleton)
%%
%% This is a SCAFFOLD. The prose between markers is for Hashem to write.
%% Claude (via /intro-help) refines wording but does not author the argument.
%%
%% Build:
%%   pdflatex introduction
%%   biber introduction
%%   pdflatex introduction
%%   pdflatex introduction
%%
%% acmart options:
%%   sigconf   — ACM SIG conference template (use this)
%%   authoryear — Harvard-style author-year citations
%%   review    — uncomment for line-numbered draft review version
%%

\documentclass[sigconf,authoryear,nonacm]{acmart}

%% nonacm=true   — removes ACM copyright block (we are not submitting to ACM)
%% sigconf       — two-column SIG conference layout
%% authoryear    — Harvard-style citations

%% ----- biblatex with Harvard styling --------------------------------------
%% acmart's authoryear option already configures natbib in author-year mode.
%% References are rendered Harvard-style by acmart's `ACM-Reference-Format`
%% bst with author-year option. If you prefer biblatex+biber, comment out
%% the natbib block and uncomment the biblatex block.

%% (Default: keep acmart's natbib-based author-year. Simpler.)

\begin{document}

\title{Beyond the Score: Investigating Prose-Level Label-Induced Bias in LLM-as-Judge Justifications}

%% TODO Hashem: replace title with your own when you finalise the RQ.

\author{Hashem}
\affiliation{
  \institution{Deakin University}
  \city{Melbourne}
  \country{Australia}}
\email{s225048421@deakin.edu.au}

\begin{abstract}
%% --- 100--150 word abstract ----------------------------------------------
%% Hashem: write last, after the body. A skeleton:
%% [1] LLM-as-judge paradigms increasingly drive evaluation of generated text.
%% [2] Saraf et al. (2025) showed numerical-score label bias.
%% [3] We extend their factorial design with a structured-justification prompt
%%     and analyse prose-level harshness independently of score.
%% [4] Pilot evidence (60 evaluations) shows X.
%% [5] Contribution statement: prose-bias detection layer; hand-coding deferred to 2.3.
%% -------------------------------------------------------------------------
TBD.
\end{abstract}

\keywords{LLM-as-judge, label-induced bias, qualitative coding, prose justification, evaluation}

\maketitle

%% =========================================================================
%% 1. BACKGROUND AND IMPORTANCE
%% =========================================================================
\section{Background and importance}
\label{sec:background}

%% Hashem: 1--2 paragraphs. Frame the NLP task. Why does LLM-as-judge matter?
%% Touch on: scale of LLM evaluation pipelines, role in benchmarking, the
%% practical stakes if judge biases propagate into model rankings.
%% Cite \citep{zheng2023judging} for the foundational paradigm.

The proliferation of large language models (LLMs) as automated evaluators
of text quality has reshaped how the field benchmarks generative systems
\citep{zheng2023judging}. \emph{[Hashem: continue here. Why does this matter
practically and scientifically? End with the framing that motivates the
paper.]}

%% =========================================================================
%% 2. CURRENT APPROACHES
%% =========================================================================
\section{Current approaches}
\label{sec:current}

%% Hashem: 1 paragraph. Brief survey of recent SOTA approaches.
%% Anchor citations: \citep{zheng2023judging,koo2024benchmarking,
%% panickssery2024llm,wataoka2024selfpreference,wang2023large}.
%% Frame: the field has built sophisticated benchmarks AND documented
%% systematic biases (self-preference, position, verbosity, etc.).

Existing work establishes that LLM judges, while approximating human
preferences in many settings \citep{zheng2023judging}, exhibit systematic
biases including self-preference \citep{panickssery2024llm,
wataoka2024selfpreference, chen2025do}, position effects
\citep{wang2023large}, and a wider catalogue of cognitive biases
\citep{koo2024benchmarking}. \emph{[Hashem: tighten and add the
LLM-judge-on-reasoning angle if you want to bring \citep{stephan2024from}
in.]}

%% =========================================================================
%% 3. IDENTIFIED LIMITATION
%% =========================================================================
\section{Identified limitation}
\label{sec:limitation}

%% Hashem: 1 paragraph. Concrete failure mode. MUST NOT be "improving accuracy".
%% The framing here is: existing label-bias work captures NUMBERS only.
%% Saraf et al. (2025) \citep{saraf2025quantifying} demonstrate numerical
%% label bias but do not analyse the prose justifications. The same
%% numerical 7/10 with savage versus gentle prose is not the same
%% evaluation. Identify the specific gap.

\citet{saraf2025quantifying} demonstrate that LLM judges exhibit pronounced
label-induced bias when scoring blog posts: the ``Claude'' label inflates
scores regardless of authorship while the ``Gemini'' label depresses them,
and false attribution can flip preference rankings by up to 50 percentage
points. Their evaluation, however, captures \emph{numerical scores only}.
\emph{[Hashem: state the precise limitation here. The same 7/10 with
savage prose vs.\ gentle prose is not the same evaluation. Prose-level
bias is invisible to score-level analysis. Be concrete.]}

%% =========================================================================
%% 4. FORMAL RESEARCH QUESTION
%% =========================================================================
\section{Formal research question}
\label{sec:rq}

%% Hashem: 1--2 sentences. The RQ is YOUR wording, not Claude's.
%% Three candidate framings to choose from (see artifacts/rq-candidates.md):
%%   (a) prose harshness shifts asymmetrically by source-model identity
%%       independent of numerical score
%%   (b) judge justifications under false-label conditions reveal
%%       rationalisation patterns
%%   (c) identity-references in justification prose track score shifts
%% Pick ONE as the headline RQ. The others can appear as secondary
%% analyses in the empirical-motivation section or methodology.

\textit{[Hashem: insert your formal RQ here. Phrase it as a falsifiable
hypothesis with a measurable outcome. State the predicted direction and
the null.]}

%% =========================================================================
%% 5. EMPIRICAL MOTIVATION
%% =========================================================================
\section{Empirical motivation}
\label{sec:motivation}

%% Hashem: 1 paragraph. This is the assignment's CORE REQUIREMENT.
%% Summarise (a) Saraf et al.'s existing score-level findings as the
%% motivating prior result, AND (b) the pilot you ran (60 evaluations
%% on 5 titles), reporting the actual prose-level shifts.
%%
%% While the pilot is not yet run, frame as "We will report the
%% pilot results from a 60-evaluation pilot run." After the pilot,
%% replace the placeholder with actual numbers.

To motivate the prose-level extension, we ran a small pilot (5 titles
$\times$ 3 judges $\times$ 4 label conditions $=$ 60 evaluations)
extending \citet{saraf2025quantifying}'s factorial design with a
structured-justification prompt that elicits a critique paragraph
alongside the numerical score---capturing a prose layer their score-only
evaluation did not. \emph{[Hashem: report pilot quantitative
result here, e.g.\ ``Mean harshness score (lexicon-derived, 1--5 ordinal)
shifted by $\Delta = X$ between true-label and false-label conditions
($p < 0.05$, paired bootstrap, $B = 10^4$), independent of the
$\Delta = Y$ point shift in numerical score. Inter-method agreement
between lexicon and Qwen-2.5-72B coder \citep{ashwin2025using,
vallejovera2025llms} was $\kappa = Z$.''] These findings indicate that
prose-level label bias exists in addition to score-level bias.}

%% =========================================================================
%% 6. CONTRIBUTION STATEMENT
%% =========================================================================
\section{Contribution statement}
\label{sec:contribution}

%% Hashem: 1 paragraph. The novel-contribution claim plus scope boundaries.
%% Include the explicit naming of Approach 3 (full hand-coding) as 2.3 HD
%% future work. This protects the contribution from over-claiming.

This paper contributes (i) the first systematic analysis of prose-level
label-induced bias in LLM-as-judge justifications, extending
\citet{saraf2025quantifying}'s numerical analysis with a
structured-justification prompt modification; (ii) a hybrid
quantitative-qualitative coding pipeline combining lexicon-based
scoring with a neutral open-weight LLM coder (Qwen-2.5-72B), explicitly
designed to avoid the circularity that would result from using
Anthropic, OpenAI, or Google models in the coder role
\citep{ashwin2025using}; and (iii) inter-method validation via
Cohen's $\kappa$ between lexicon and LLM-coder dimensions. We
explicitly defer human hand-coding validation to future work in the
companion High Distinction task as a complementary trust-building
layer. \emph{[Hashem: tighten this. The contribution should be
falsifiable, narrowly scoped, and honest about what's not being
claimed.]}

%% =========================================================================
%% Bibliography
%% =========================================================================
\bibliographystyle{ACM-Reference-Format}
\bibliography{references}

\end{document}
